% GENERATED FILE. Edit canonical lessons, then run npm run generate:books.
% canonical-source-hash: fnv1a64:8c86b8001d8f6692
% canonical-lessons: ML-C275-borati, ML-C275-mula3, ML-C275-catavu, ML-C275-mannappittam, ML-C275-raktasammarddam

\chapter{Boredom, Lump, Bruise, Jaundice, Blood pressure}
\label{ch:ml-boredom-lump-bruise-jaundice-blood-pressure}
\hlchaptermodality{275}

\begin{chapteropening}
\textbf{By the end of this chapter:} \emph{I can say boredom, a lump, a bruise, jaundice and blood pressure.}

Complete the last lesson of chapter 275: I can say boredom, a lump, a bruise, jaundice and blood pressure.
\end{chapteropening}

\section[\texorpdfstring{\ml{ബോറടി} (bōraṭi)}{ബോറടി (bōraṭi)}]{\ml{ബോറടി} (bōraṭi) — boredom}

\label{lesson:ML-C275-borati}



\begin{quote}
Before the new one: say the Malayalam for a burn, then the Malayalam for diabetes.
\end{quote}

\subsection*{You'll want to know: \ml{ബോറടി}}
\textbf{\ml{ബോറടി}} — \emph{bōraṭi} — \textquotedblleft{}boredom\textquotedblright{}.

\begin{grammarlens}[title={What you've built}]
One more word for this chapter.
\end{grammarlens}

\subsection*{Your turn}
\begin{itemize}
\raggedright
  \item \emph{Say it:} \emph{bōraṭi}
  \item \emph{Say it:} \emph{bōraṭi}, once more
  \item \emph{Say it:} say \emph{bōraṭi}
  \item \emph{Recall:} say the Malayalam for a burn, then the Malayalam for diabetes, then say \emph{bōraṭi} again
\end{itemize}

\begin{tcolorbox}[breakable,colback=teal!4,colframe=teal!35!black,title={Before you move on}]
What does \ml{ബോറടി} mean? (\textquotedblleft{}Boredom\textquotedblright{}.) Say it once more.
\end{tcolorbox}
\section[\texorpdfstring{\ml{മുഴ} (muḻa)}{മുഴ (muḻa)}]{\ml{മുഴ} (muḻa) — a lump, a swelling bump}

\label{lesson:ML-C275-mula3}



\begin{quote}
Before the new one: say the Malayalam for diabetes, then the Malayalam for boredom.
\end{quote}

\subsection*{You'll want to know: \ml{മുഴ}}
\textbf{\ml{മുഴ}} — \emph{muḻa} — \textquotedblleft{}a lump, a swelling bump\textquotedblright{}.

\begin{grammarlens}[title={What you've built}]
One more word for this chapter.
\end{grammarlens}

\subsection*{Your turn}
\begin{itemize}
\raggedright
  \item \emph{Say it:} \emph{muḻa}
  \item \emph{Say it:} \emph{muḻa}, once more
  \item \emph{Say it:} say \emph{muḻa}
  \item \emph{Recall:} say the Malayalam for diabetes, then the Malayalam for boredom, then say \emph{muḻa} again
\end{itemize}

\begin{tcolorbox}[breakable,colback=teal!4,colframe=teal!35!black,title={Before you move on}]
What does \ml{മുഴ} mean? (\textquotedblleft{}A lump, a swelling bump\textquotedblright{}.) Say it once more.
\end{tcolorbox}
\section[\texorpdfstring{\ml{ചതവ്} (catavŭ)}{ചതവ് (catavŭ)}]{\ml{ചതവ്} (catavŭ) — a bruise}

\label{lesson:ML-C275-catavu}



\begin{quote}
Before the new one: say the Malayalam for boredom, then the Malayalam for a lump.
\end{quote}

\subsection*{You'll want to know: \ml{ചതവ്}}
\textbf{\ml{ചതവ്}} — \emph{catavŭ} — \textquotedblleft{}a bruise\textquotedblright{}.

\begin{grammarlens}[title={What you've built}]
One more word for this chapter.
\end{grammarlens}

\subsection*{Your turn}
\begin{itemize}
\raggedright
  \item \emph{Say it:} \emph{catavŭ}
  \item \emph{Say it:} \emph{catavŭ}, once more
  \item \emph{Say it:} say \emph{catavŭ}
  \item \emph{Recall:} say the Malayalam for boredom, then the Malayalam for a lump, then say \emph{catavŭ} again
\end{itemize}

\begin{tcolorbox}[breakable,colback=teal!4,colframe=teal!35!black,title={Before you move on}]
What does \ml{ചതവ്} mean? (\textquotedblleft{}A bruise\textquotedblright{}.) Say it once more.
\end{tcolorbox}
\section[\texorpdfstring{\ml{മഞ്ഞപ്പിത്തം} (maññappittaṁ)}{മഞ്ഞപ്പിത്തം (maññappittaṁ)}]{\ml{മഞ്ഞപ്പിത്തം} (maññappittaṁ) — jaundice}

\label{lesson:ML-C275-mannappittam}



\begin{quote}
Before the new one: say the Malayalam for a lump, then the Malayalam for a bruise.
\end{quote}

\subsection*{You'll want to know: \ml{മഞ്ഞപ്പിത്തം}}
\textbf{\ml{മഞ്ഞപ്പിത്തം}} — \emph{maññappittaṁ} — \textquotedblleft{}jaundice\textquotedblright{}.

\begin{grammarlens}[title={What you've built}]
One more word for this chapter.
\end{grammarlens}

\subsection*{Your turn}
\begin{itemize}
\raggedright
  \item \emph{Say it:} \emph{maññappittaṁ}
  \item \emph{Say it:} \emph{maññappittaṁ}, once more
  \item \emph{Say it:} say \emph{maññappittaṁ}
  \item \emph{Recall:} say the Malayalam for a lump, then the Malayalam for a bruise, then say \emph{maññappittaṁ} again
\end{itemize}

\begin{tcolorbox}[breakable,colback=teal!4,colframe=teal!35!black,title={Before you move on}]
What does \ml{മഞ്ഞപ്പിത്തം} mean? (\textquotedblleft{}Jaundice\textquotedblright{}.) Say it once more.
\end{tcolorbox}
\section[\texorpdfstring{\ml{രക്തസമ്മർദ്ദം} (raktasammarddaṁ)}{രക്തസമ്മർദ്ദം (raktasammarddaṁ)}]{\ml{രക്തസമ്മർദ്ദം} (raktasammarddaṁ) — blood pressure}

\label{lesson:ML-C275-raktasammarddam}



\begin{quote}
Before the new one: say the Malayalam for a bruise, then the Malayalam for jaundice.
\end{quote}

\subsection*{You'll want to know: \ml{രക്തസമ്മർദ്ദം}}
\textbf{\ml{രക്തസമ്മർദ്ദം}} — \emph{raktasammarddaṁ} — \textquotedblleft{}blood pressure\textquotedblright{}.

\begin{grammarlens}[title={What you've built}]
One more word for this chapter.
\end{grammarlens}

\subsection*{Your turn}
\begin{itemize}
\raggedright
  \item \emph{Say it:} \emph{raktasammarddaṁ}
  \item \emph{Say it:} \emph{raktasammarddaṁ}, once more
  \item \emph{Say it:} say \emph{raktasammarddaṁ}
  \item \emph{Recall:} say the Malayalam for a bruise, then the Malayalam for jaundice, then say \emph{raktasammarddaṁ} again
\end{itemize}

\begin{tcolorbox}[breakable,colback=teal!4,colframe=teal!35!black,title={Before you move on}]
What does \ml{രക്തസമ്മർദ്ദം} mean? (\textquotedblleft{}Blood pressure\textquotedblright{}.) Say it once more.
\end{tcolorbox}
